\section{Related Work}

\paragraph{Exact tensor network contraction.}
Classical work on exact tensor network contraction has concentrated on contraction ordering, contraction trees, treewidth-aware planning, and slicing strategies. A consistent takeaway from this line of research is that the dominant difficulty is usually not the original tensors themselves, but the size of the intermediate states they induce and the resulting time and memory complexity~\cite{markov2008simulating,dumitrescu2018benchmarking,jakes2019carving,gray2021hyperoptimized}. Index slicing further shows that when a monolithic contraction is too expensive, one can cut selected degrees of freedom and solve many smaller subproblems in parallel or in batches~\cite{huang2021efficient}. These ideas directly motivate our setting: separator growth is the true bottleneck, and blockwise organization can be an important systems-level tool when full materialization is required. However, the primary objective of those methods is still exact contraction rather than approximate acceleration under an unchanged dense-output requirement.

\paragraph{Approximate contraction and intermediate-state compression.}
When exact contraction becomes prohibitive, approximate tensor-network methods typically manage complexity by compressing intermediate objects. On regular lattices, work such as tensor renormalization group (TRG) established that effective truncation of otherwise intractable intermediate states is essential to scalable approximate contraction~\cite{levin2007trg}. More recently, these ideas have been extended from regular topologies to general graphs. Pan et al.\ proposed a general-purpose approximate contraction algorithm for arbitrary connected structures~\cite{pan2020contracting}; Gray and Chan incorporated approximation directly into contraction planning so that compression and contraction order can be co-optimized~\cite{gray2024hyperoptimized}; and Ma et al.\ developed a density-matrix-style approximation framework for general graphs~\cite{ma2024approximate}. These methods show that approximation on irregular graphs can be treated systematically. Our focus differs in that the final output tensor must still be fully produced, so the natural place for approximation is the internal separator state rather than the output boundary.

\paragraph{Randomized low-rank approximation and sketching.}
Our implementation also draws on randomized numerical linear algebra. Randomized low-rank approximation demonstrates that dominant subspaces can often be captured efficiently through random projections, followed by stable decomposition in a much smaller projected space, without first computing an expensive exact factorization~\cite{halko2011finding}. More broadly, Count Sketch and Tensor Sketch illustrate how randomized linear maps can compress high-dimensional objects while preserving useful structure~\cite{charikar2004finding,pham2013fast}. Very recent work has begun to discuss sketching directly inside tensor network contraction~\cite{heddes2026approximating}. In our setting, randomized compression is an implementation tool inside a separator-state framework; it does not change the problem definition.

\paragraph{Connections to elimination in graphical models.}
At the algorithmic level, our approach is also related to bucket elimination and mini-bucket methods in graphical models~\cite{dechter1996bucket,dechter1997mini}. Those methods emphasize that what must be propagated upward is not every local detail, but only the interface information that remains relevant to the rest of the graph. When the interface becomes too large, one needs a controlled approximation at the intermediate level. We revisit this viewpoint in the setting of open-leg tensor networks with full dense materialization, and we explicitly treat preservation of the output boundary as part of the method definition.
